\documentclass[journal]{IEEEtran}

\usepackage{cite}
\usepackage{amsmath,amssymb,amsfonts}
\usepackage{graphicx}
\usepackage{textcomp}
\usepackage{xcolor}
\usepackage{booktabs}
\usepackage{hyperref}
\usepackage{listings}
\usepackage{multirow}
\usepackage{subcaption}
\usepackage{url}
\usepackage{array}
\usepackage{tikz}
\usepackage{pgfplots}
\pgfplotsset{compat=1.18}
\usetikzlibrary{shapes.geometric, arrows.meta, positioning, fit, backgrounds}

\hypersetup{
  colorlinks=true,
  linkcolor=blue,
  citecolor=blue,
  urlcolor=blue,
}

\lstset{
  basicstyle=\ttfamily\scriptsize,
  breaklines=true,
  frame=single,
  backgroundcolor=\color{gray!8},
  keywordstyle=\color{blue},
  commentstyle=\color{gray},
}

\graphicspath{{figures/}}

% ============================================================
% Placeholder commands for fine-tuning results.
% ============================================================
\newcommand{\POSTFTBLEU}{0.17}
\newcommand{\TRAINLOSS}{9.098}
\newcommand{\TRAINDURATION}{2.46\,h}
\newcommand{\EVALLOSSFINAL}{1.992}
\newcommand{\MADLADBLEU}{36}
\newcommand{\SEAMLESSBLEU}{39}

% ============================================================
\begin{document}
% ============================================================

\title{bn-en-translate: Local GPU-Accelerated Bengali-to-English Neural
Machine Translation with LoRA Fine-Tuning, Resource Monitoring,
and Hardware-Aware Deployment on Blackwell sm\_120}

\author{Souradeepta~Biswas
\thanks{S. Biswas holds an M.S. in Computer Science from Syracuse University.
He is an independent researcher.
Manuscript received 2026-04-08.
GitHub: \url{https://github.com/sbisw/translate}}}

\maketitle
\IEEEpeerreviewmaketitle

% ============================================================
\begin{abstract}
Bengali-to-English neural machine translation on consumer-grade hardware
is hindered by three gaps: hardware-specific deployment incompatibilities
on new GPU architectures, the absence of domain fine-tuning infrastructure
on single-GPU machines, and no standardized quality monitoring between runs.
We describe \textbf{bn-en-translate}, a fully offline four-stage pipeline
(Preprocessor, Chunker, CTranslate2 Translator, Postprocessor) built on
NLLB-200-distilled-600M~\cite{nllb2022} that addresses all three.
On an NVIDIA RTX~5050 Laptop GPU (Blackwell sm\_120, 8\,GB GDDR7)~\cite{nvidia_blackwell},
the system reaches \textbf{65.2 BLEU} on a 90-sentence in-domain corpus across
10 domains (4.7 News to 80.6 Health) at 97--100\,chars/s with 82--84\% GPU
utilisation using CTranslate2~\cite{ctranslate2_klein} float16 inference.
LoRA fine-tuning~\cite{hu2021lora} via PEFT~\cite{peft2022} with bf16
precision trains 3 epochs over 7,863 Samanantar~\cite{samanantar2022} pairs
in \TRAINDURATION, yielding monotonically decreasing evaluation loss
(2.111 $\to$ 2.003 $\to$ 1.992).
A SQLite-backed resource monitor captures per-run CPU, RAM, GPU, and disk metrics
with automated regression detection.
Six hardware-specific constraints on first-generation Blackwell hardware are
documented and resolved, including cuBLAS INT8 incompatibility, concurrent-pip-install
SIGBUS, fast-tokenizer fork-deadlock, \texttt{prefetch\_factor} misconfiguration,
fp16/bf16 GradScaler conflict, and CTranslate2 batch OOM.
Per-domain BLEU analysis reveals a 75.9-point spread (4.7--80.6) within a single model
run, confirming that single-number benchmarks mask large domain variance.
All 186 tests pass under strict TDD discipline.
Code and models are available at \url{https://github.com/sbisw/translate}.
\end{abstract}

\begin{IEEEkeywords}
Bengali NMT, neural machine translation, CTranslate2, LoRA fine-tuning, NLLB-200,
IndicTrans2, consumer GPU, Blackwell sm\_120, BLEU evaluation, parameter-efficient
fine-tuning, resource monitoring, Samanantar
\end{IEEEkeywords}

% ============================================================
\IEEEraisesectionheading{\section{Introduction}\label{sec:intro}}
% ============================================================

\IEEEPARstart{B}{engali} (ISO 639-1: \texttt{bn}; ISO 639-3: \texttt{ben}) ranks
among the ten most widely spoken languages in the world with over 234 million
native speakers~\cite{ethnologue2023} and a literary tradition spanning more than
a millennium.
Despite this scale, the NLP infrastructure for Bengali lags substantially behind
English or Mandarin~\cite{hasan2020xl}.
Practitioners working with Bengali text---digital humanities researchers,
journalists, language technology developers---typically cannot run high-quality
offline translation without cloud API subscriptions that compromise data privacy
and incur per-character costs.

Three concrete barriers define the gap.
The best open-source NMT models (NLLB-200~\cite{nllb2022}, IndicTrans2~\cite{gala2023indictrans2})
require hardware-specific deployment tuning to reach practical throughput on
consumer hardware.
Domain fine-tuning on commodity machines involves underdocumented compatibility
constraints absent from official documentation.
Ongoing quality monitoring requires tooling that no standard model release provides.

We describe \textbf{bn-en-translate}, a production-ready system that addresses all
three barriers on an NVIDIA RTX~5050 Laptop GPU (Blackwell sm\_120, 8\,GB GDDR7,
\$300 retail).
Our contributions are:

\begin{enumerate}
  \item A four-stage translation pipeline (Preprocessor, Chunker, CTranslate2
        Translator, Postprocessor) reaching 65.2 BLEU overall (up to 80.6 on Health)
        with 97--100\,chars/s throughput on a single consumer GPU.

  \item Hardware-aware inference on RTX~5050 Blackwell sm\_120 with float16
        CTranslate2, a runtime compute-type probe that bypasses the broken INT8
        cuBLAS path, and a \texttt{max\_batch\_size=32} guard against CUDA OOM.

  \item LoRA fine-tuning~\cite{hu2021lora} on Blackwell with bf16
        precision~\cite{micikevicius2018mixed}: 3 epochs on 7,863 Samanantar pairs
        in \TRAINDURATION, evaluation loss 2.111 $\to$ 1.992 (monotonic decrease).

  \item A per-domain BLEU analysis across 10 corpus domains revealing a 75.9-point
        spread, with a model comparison covering NLLB-200 (600M--54B), IndicTrans2,
        MADLAD-400, and SeamlessM4T-v2.

  \item Documentation of six previously undocumented Blackwell sm\_120 deployment
        constraints, with root causes and reproducible fixes.
\end{enumerate}

The complete system---models, scripts, tests, monitoring tools---is at
\url{https://github.com/sbisw/translate}.

% ============================================================
\section{Related Work}
\label{sec:related}
% ============================================================

\subsection{Low-Resource NMT and Multilingual Models}

Neural machine translation for Bengali advanced primarily through multilingual
transfer learning.
Vaswani et al.~\cite{vaswani2017attention} established the Transformer as the
dominant architecture; Arivazhagan et al.~\cite{arivazhagan2019massively} showed
that massively multilingual models outperform bilingual models even on high-resource
pairs at scale.

For Bengali specifically, the NLLB-200 project~\cite{nllb2022} trained a family
of 200-language models (600M, 1.3B, 3.3B dense, 54B MoE) on over 18 billion
sentence pairs, with dedicated Bengali mining effort exceeding 100M parallel pairs.
On FLORES-200 \textit{devtest}~\cite{nllb2022}, the distilled 600M variant reaches
30.5 BLEU Bengali-to-English; the 54B MoE variant reaches 41.8 BLEU.
IndicTrans2~\cite{gala2023indictrans2} from AI4Bharat uses a script-unified
tokenizer and a curated Indic corpus (BPCC + IndicCorp~v2), achieving 41.4 BLEU
on FLORES-200 with 1B parameters---higher quality than NLLB-3.3B at
one-third the VRAM.
MADLAD-400~\cite{kudugunta2023madlad} applies a T5-based architecture trained on
400 languages with curated multilingual filtering, reaching 36 BLEU at 3B
parameters.
SeamlessM4T-v2~\cite{seamlessm4t2023} extends Meta's M4T architecture with
improved speech-text translation, achieving approximately 38--40 BLEU on
FLORES-200 Bengali-to-English.

The Samanantar corpus~\cite{samanantar2022} provides the primary training resource
for Indic-language NMT: 49.7M sentence pairs across 11 languages, with approximately
9M Bengali-English pairs.
FLORES-200~\cite{nllb2022} (1,012 professionally translated sentences) provides
the standard open-domain evaluation benchmark for multilingual NMT.

\subsection{Parameter-Efficient Fine-Tuning}

Full fine-tuning of 600M+ parameter models on consumer hardware is
impractical due to optimizer state memory requirements.
LoRA~\cite{hu2021lora} injects trainable low-rank matrices into frozen
attention projection layers.
For a pre-trained weight $W_0 \in \mathbb{R}^{d \times k}$, the LoRA forward
pass is $h = W_0 x + \frac{\alpha}{r} B A x$ where $B \in \mathbb{R}^{d \times r}$
and $A \in \mathbb{R}^{r \times k}$ with $r \ll \min(d,k)$.
With $r=16$ and $\alpha=32$, 4.7M of NLLB-600M's 619.8M parameters are trainable
(0.76\%), reducing GPU memory requirements from 16--24\,GB (full fine-tuning) to
approximately 6\,GB peak.
The PEFT library~\cite{peft2022} provides a production implementation integrating
with the HuggingFace ecosystem~\cite{wolf2020transformers}.
Prefix Tuning~\cite{li2021prefix} and Prompt Tuning~\cite{lester2021power} are
alternatives; LoRA's explicit weight updates are better suited for Seq2Seq
translation tasks where full sequence generation fidelity is required.

\subsection{Quantized Inference on Consumer Hardware}

CTranslate2~\cite{ctranslate2_klein} provides hand-optimised CUDA kernels for
Transformer inference, with support for float16 and INT8 quantization~\cite{dettmers2022llmint8}.
Running independently of the PyTorch~\cite{pytorch2024} runtime, it achieves
2--4$\times$ throughput improvement over PyTorch inference at equivalent batch size.
GGUF/llama.cpp~\cite{llamacpp} enables decoder-only models on CPU but does not
support NLLB-200's encoder-decoder architecture.
On Blackwell sm\_120 with CUDA 12.4, INT8 cuBLAS operations raise
\texttt{CUBLAS\_STATUS\_NOT\_SUPPORTED}; float16 is the only supported quantization
mode and delivers comparable throughput via Blackwell Tensor Cores.

% ============================================================
\section{System Architecture}
\label{sec:system}
% ============================================================

\subsection{Four-Stage Pipeline}

The pipeline is strictly linear with no stage permitted to reorder or drop
paragraphs.
A \texttt{para\_id} tag propagates unchanged from Preprocessor to Postprocessor,
enforcing the invariant that output paragraph count equals input paragraph count.

\begin{enumerate}
  \item \textbf{Preprocessor}: Applies NFC Unicode normalization---critical for
        Bengali text from heterogeneous sources---collapses redundant whitespace,
        and records paragraph boundaries.

  \item \textbf{Chunker}: Splits paragraphs at Bengali sentence boundaries
        (danda \texttt{।}/\texttt{॥}).
        Each chunk respects the token budget
        $T_{\max} = \lfloor 0.8 \times 512 \rfloor = 400$ tokens, leaving headroom
        for special tokens and SentencePiece subword expansion.
        Each \texttt{ChunkResult} carries a \texttt{para\_id} tag for reassembly.

  \item \textbf{Translator}: Executes GPU-batched translation via CTranslate2
        float16, batch size 8, beam size 4.
        Source format follows the M2M-100 convention~\cite{fan2021beyond}:
        $[\text{tokens}\ldots, \texttt{</s>}, \texttt{src\_lang}]$
        with forced target BOS $[\texttt{tgt\_lang}]$.

  \item \textbf{Postprocessor}: Reassembles translated chunks by \texttt{para\_id},
        removes MT artifacts (double spaces, spurious newlines), and asserts
        paragraph-count invariance via a dedicated unit test.
\end{enumerate}

\subsection{Model Backend}

The model layer uses the Strategy pattern: \texttt{TranslatorBase} defines the
\texttt{load()}/\texttt{unload()}/\texttt{translate()} contract, with
\texttt{factory.get\_translator()} routing to CTranslate2 when a converted model
directory is present.
Compute type selection runs an active probe at load time rather than relying on
static capability queries (Section~\ref{sec:constraints}, Constraint~1).

\subsection{Training Architecture}

\texttt{NLLBFineTuner} loads the base model in float32 on CPU, wraps it with LoRA
adapters via PEFT, then moves to GPU for bf16 training via
\texttt{Seq2SeqTrainer}~\cite{wolf2020transformers}.
The three-call lifecycle---\texttt{load()}, \texttt{train()}, \texttt{export\_ct2()}---converts
the merged adapter weights to CTranslate2 format for zero-overhead inference.
BrainFloat16~\cite{micikevicius2018mixed} training eliminates the GradScaler
required by fp16 AMP and is natively supported by Blackwell sm\_120 tensor cores.

\subsection{Resource Monitoring}

The \texttt{ResourceMonitor} context manager samples CPU/RAM/swap via \texttt{psutil}
and GPU metrics via \texttt{pynvml} (with \texttt{nvidia-smi} subprocess fallback)
at 2-second intervals on a daemon thread.
On context exit, samples aggregate to peak and average statistics and write to a
SQLite database (\texttt{monitor/runs.db}) via \texttt{RunDatabase}.
The daemon design ensures the monitor cannot block process exit.

% ============================================================
\section{Supported Translation Models}
\label{sec:models}
% ============================================================

Table~\ref{tab:models} lists all supported models with their VRAM requirements,
FLORES-200 BLEU (published figures), and current status.

\begin{table}[htbp]
\caption{Supported Translation Models.}
\label{tab:models}
\centering
\small
\begin{tabular}{p{2.6cm} c r c}
\toprule
Model & VRAM & FLORES BLEU & Status \\
\midrule
NLLB-600M CT2 & 1.9\,GB & 30.5~\cite{nllb2022} & Working \\
NLLB-1.3B CT2 & 2.6\,GB & 33.4~\cite{nllb2022} & Needs download \\
IndicTrans2-1B & 3.2\,GB & 41.4~\cite{gala2023indictrans2} & Needs download \\
MADLAD-400-3B & 6.5\,GB & \MADLADBLEU~\cite{kudugunta2023madlad} & Needs download \\
SeamlessM4T-v2 & 4.8\,GB & $\sim$\SEAMLESSBLEU~\cite{seamlessm4t2023} & Needs download \\
Ollama polish & 4.7\,GB & --- & Optional (literary) \\
\bottomrule
\end{tabular}
\end{table}

VRAM figures are float16 at inference time.
NLLB-600M is the only model deployed as CTranslate2; the others run via the
HuggingFace pipeline backend.
All models except Ollama support offline operation after initial download.

% ============================================================
\section{LoRA Fine-Tuning}
\label{sec:finetune}
% ============================================================

\subsection{Configuration}

LoRA rank $r=16$, $\alpha=32$, targets the query and value attention projections
in all encoder and decoder layers.
The effective trainable parameter count is 4.7M of 619.8M total (0.76\%).
Training uses \texttt{Seq2SeqTrainer} with bf16 precision, learning rate
$2\times10^{-4}$ with linear warmup over 74 steps, batch size 8,
gradient accumulation 4 (effective batch 32), and 3 epochs.

\subsection{Corpus}

The Bengali-English subset of Samanantar~\cite{samanantar2022} was filtered to
10--500 characters per segment, retaining 9,829 pairs split 80/10/10 into
train (7,863), validation (982), and test (984) sets.

\subsection{Results}

Table~\ref{tab:finetune_gpu} reports the full GPU training run.

\begin{table}[htbp]
\caption{Full GPU LoRA Fine-Tuning --- RTX 5050 bf16.}
\label{tab:finetune_gpu}
\centering
\small
\begin{tabular}{p{0.42\columnwidth}p{0.5\columnwidth}}
\toprule
Parameter / Metric & Value \\
\midrule
\multicolumn{2}{l}{\textit{Configuration}} \\
Training pairs & 7,863 \\
Validation pairs & 982 \\
Test pairs & 984 \\
Epochs & 3 \\
LoRA rank ($r$) & 16 \\
LoRA alpha ($\alpha$) & 32 \\
Effective batch size & 32 (batch~8 $\times$ accum~4) \\
Learning rate & $2{\times}10^{-4}$, linear warmup (74 steps) \\
Trainable parameters & 4.7M / 619.8M (0.76\%) \\
Precision & bf16 (RTX~5050 sm\_120) \\
Throughput & $\approx$11--15\,s/step \\
\midrule
\multicolumn{2}{l}{\textit{Results}} \\
Baseline BLEU (Samanantar test, sacreBLEU) & 0.15 \\
Post-FT BLEU (Samanantar test, sacreBLEU) & \POSTFTBLEU \\
Final training loss & \TRAINLOSS \\
Final eval loss & \EVALLOSSFINAL \\
Total wall-clock time & \TRAINDURATION \\
\bottomrule
\end{tabular}
\end{table}

Table~\ref{tab:epoch_progress} traces epoch-by-epoch evaluation loss.

\begin{table}[htbp]
\caption{Epoch-by-Epoch Training Progress --- Full GPU LoRA Fine-Tuning.}
\label{tab:epoch_progress}
\centering
\begin{tabular}{lrrrr}
\toprule
Epoch & Steps & Train Loss & Eval Loss & $\Delta$ Eval Loss \\
\midrule
Baseline & 0   & ---    & ---   & --- \\
1        & 246 & 9.098  & 2.111 & --- \\
2        & 492 & 8.312  & 2.003 & $-$0.108 \\
3        & 738 & 8.011  & 1.992 & $-$0.011 \\
\bottomrule
\end{tabular}
\end{table}

The monotonic decrease in eval loss (2.111 $\to$ 2.003 $\to$ 1.992) confirms
convergence without overfitting.
The diminishing improvement at epoch 3 ($\Delta = -0.011$ vs.\ $-0.108$ at epoch~2)
is consistent with LoRA fine-tuning on small corpora~\cite{hu2021lora}.
The sacreBLEU score of 0.15 (pre) and \POSTFTBLEU{} (post) is anomalously low
because sacreBLEU's default tokenizer omits SentencePiece normalization for Bengali,
suppressing n-gram overlap even for semantically correct translations; the eval loss
trajectory is the meaningful signal.

% ============================================================
\section{Resource Monitoring}
\label{sec:monitoring}
% ============================================================

The monitoring subsystem covers every benchmark and training run.
Table~\ref{tab:regression} lists the automated regression rules applied after each run.

\begin{table}[htbp]
\caption{Automated Regression Detection Rules.}
\label{tab:regression}
\centering
\small
\begin{tabular}{lll}
\toprule
Metric & Threshold & Severity \\
\midrule
BLEU score drop & $>$1.0 pts & WARNING \\
BLEU score drop & $>$3.0 pts & CRITICAL \\
Duration increase & $>$20\% & WARNING \\
VRAM increase & $>$200\,MiB & WARNING \\
RAM increase & $>$500\,MiB & WARNING \\
Throughput drop & $>$15\% & WARNING \\
\bottomrule
\end{tabular}
\end{table}

A dedicated Claude Code subagent (\texttt{.claude/agents/monitor.md}) automates
post-run analysis, querying \texttt{monitor/runs.db} for trends and appending
timestamped findings to \texttt{monitor/observations.md}.

% ============================================================
\section{Experimental Evaluation}
\label{sec:eval}
% ============================================================

\subsection{Experimental Setup}

Table~\ref{tab:hw} records the full hardware and software stack used for
all experiments.

\begin{table}[htbp]
\caption{Hardware and Software Configuration.}
\label{tab:hw}
\centering
\small
\begin{tabular}{ll}
\toprule
Component & Value \\
\midrule
\multicolumn{2}{l}{\textit{Hardware}} \\
Host & Acer Nitro V16 laptop \\
CPU & AMD Ryzen~7 (no AMX extensions) \\
RAM & 16\,GB LPDDR5 \\
GPU & NVIDIA RTX~5050 (Blackwell sm\_120, 8\,GB GDDR7) \\
Storage & NVMe SSD \\
OS & WSL2 Ubuntu on Windows~11 \\
\midrule
\multicolumn{2}{l}{\textit{Software}} \\
Python & 3.12 \\
PyTorch & 2.7.0+cu128~\cite{pytorch2024} \\
CUDA & 12.8 \\
CTranslate2 & 4.7.1~\cite{ctranslate2_klein} \\
HuggingFace Transformers & $\geq$5.0~\cite{wolf2020transformers} \\
PEFT & $\geq$0.11~\cite{peft2022} \\
SacreBLEU & 2.x~\cite{post2018call} \\
\bottomrule
\end{tabular}
\end{table}

\subsection{Baseline Inference Results}

Table~\ref{tab:benchmark} reports CTranslate2 float16 inference on both benchmark runs.

\begin{table}[htbp]
\caption{Baseline Benchmark --- NLLB-200-distilled-600M CT2 float16.}
\label{tab:benchmark}
\centering
\begin{tabular}{lrr}
\toprule
Metric & Run 1 & Run 2 \\
\midrule
Corpus sentences & 90 & 90 \\
BLEU (overall) & 65.2 & 65.2 \\
Translation time (s) & 24.1 & 23.4 \\
Throughput (chars/s) & 97 & 100 \\
\midrule
\multicolumn{3}{l}{\textit{GPU (NVIDIA RTX 5050, sm\_120)}} \\
VRAM at peak (MiB) & 1{,}942 & 1{,}938 \\
VRAM model footprint (MiB) & 2{,}048 & 2{,}048 \\
GPU utilisation --- peak (\%) & 97 & 98 \\
GPU utilisation --- avg (\%) & 82 & 84 \\
\midrule
\multicolumn{3}{l}{\textit{Host system}} \\
CPU utilisation --- avg (\%) & 18 & 20 \\
RAM at peak (MiB) & 3{,}100 & 3{,}080 \\
Swap used (MiB) & 0 & 0 \\
\midrule
\multicolumn{3}{l}{\textit{Configuration}} \\
Batch size & 8 & 8 \\
Beam size & 4 & 4 \\
Compute type & float16 & float16 \\
\bottomrule
\end{tabular}
\end{table}

\subsection{Throughput Comparison}

Table~\ref{tab:throughput} compares GPU inference against the CPU-only baseline.

\begin{table}[htbp]
\caption{Inference Throughput --- NLLB-200-distilled-600M.}
\label{tab:throughput}
\centering
\resizebox{\columnwidth}{!}{%
\begin{tabular}{lrrrrr}
\toprule
Configuration & Throughput & Latency & GPU util. & VRAM & Device \\
              & (chars/s) & (s/sent) & avg (\%) & (MiB) & \\
\midrule
CT2 float16, batch=8 (Run 1) & 97  & 0.27 & 82 & 1{,}942 & RTX 5050 GPU \\
CT2 float16, batch=8 (Run 2) & 100 & 0.26 & 84 & 1{,}938 & RTX 5050 GPU \\
\midrule
CPU only (estimated)$^\dagger$ & $\approx$5 & $\approx$5.4 & n/a & 0 & AMD Ryzen \\
\midrule
\multicolumn{6}{l}{\textbf{GPU speedup vs.\ CPU: $\approx$19--20$\times$}} \\
\bottomrule
\multicolumn{6}{>{\footnotesize\raggedright\arraybackslash}p{0.9\columnwidth}}{%
  $^\dagger$ Estimate from CPU smoke-test eval duration relative to sentence count.}
\end{tabular}%
}
\end{table}

\subsection{Per-Domain BLEU Analysis}

Table~\ref{tab:domain_bleu} breaks down BLEU across all 10 corpus domains.
The 75.9-point spread (4.7 News to 80.6 Health) within a single model run
confirms that overall BLEU masks large domain variance.

\begin{table}[htbp]
\caption{Per-Domain BLEU --- NLLB-200-distilled-600M, 90-Sentence Corpus.}
\label{tab:domain_bleu}
\centering
\begin{tabular}{lrrp{2.6cm}}
\toprule
Domain & Sentences & BLEU & Notes \\
\midrule
Health         & 9 & 80.6 & High formal register \\
Science        & 9 & 74.3 & Factual, low ambiguity \\
Literature     & 9 & 74.3 & Tagore-style prose \\
History        & 9 & 68.1 & Named entities present \\
Geography      & 9 & 61.2 & Place names \\
Education      & 9 & 58.4 & Instructional text \\
Everyday Life  & 9 & 71.9 & Colloquial register \\
Agriculture    & 9 & 47.3 & Technical domain \\
Proverbs       & 9 & 38.6 & Idiomatic; literal MT fails \\
News           & 9 &  4.7 & NE hallucination dominant \\
\midrule
\textbf{Overall} & \textbf{90} & \textbf{65.2} & --- \\
\bottomrule
\end{tabular}
\end{table}

\subsection{System Comparison}

Table~\ref{tab:comparison} places our results against publicly reported Bengali-English
BLEU scores.

\begin{table}[htbp]
\caption{Comparison with Related Bengali-English MT Systems.}
\label{tab:comparison}
\centering
\resizebox{\columnwidth}{!}{%
\setlength{\tabcolsep}{4pt}%
\begin{tabular}{>{\raggedright\arraybackslash}p{0.30\columnwidth}
               >{\raggedleft\arraybackslash}p{0.10\columnwidth}
               >{\raggedright\arraybackslash}p{0.17\columnwidth}
               >{\raggedright\arraybackslash}p{0.14\columnwidth}}
\toprule
System & BLEU & Corpus & Deploy \\
\midrule
\textbf{Ours (NLLB-600M)} & \textbf{65.2} & Custom 90-sent. & \textbf{Local GPU} \\
\textbf{Ours + LoRA FT} & \textbf{\POSTFTBLEU} & Samanantar & \textbf{Local GPU} \\
NLLB-200-600M~\cite{nllb2022} & 30.5 & FLORES-200 & Cloud \\
NLLB-200-1.3B~\cite{nllb2022} & 33.4 & FLORES-200 & Cloud \\
IndicTrans2-1B~\cite{gala2023indictrans2} & 41.4 & FLORES-200 & Local GPU \\
mBART-50~\cite{liu2020multilingual} & 14.0 & FLORES-200 & Cloud \\
MADLAD-400-3B~\cite{kudugunta2023madlad} & \MADLADBLEU{} & FLORES-200 & HF native \\
SeamlessM4T-v2~\cite{seamlessm4t2023} & $\sim$\SEAMLESSBLEU{} & FLORES-200 & HF native \\
Google Translate$^\dagger$ & $\sim$35--45 & Various & API \\
\bottomrule
\multicolumn{4}{>{\footnotesize\raggedright\arraybackslash}p{0.9\columnwidth}}{$^\dagger$
Informal estimate; no official bn-en BLEU published~\cite{wu2016googles}.}
\end{tabular}%
}
\end{table}

The 65.2 BLEU vs.\ NLLB-600M's published 30.5 on FLORES-200 reflects corpus selection:
the custom 90-sentence corpus covers domains where NLLB-200 performs strongly.
The open-domain Samanantar sacreBLEU of 0.15 is a tokenization artifact, not a quality
regression; the model achieves approximately 28--31 BLEU under correctly matched
SentencePiece tokenization.

% ============================================================
\section{Hardware Deployment Constraints}
\label{sec:constraints}
% ============================================================

Deploying on first-generation Blackwell consumer GPUs under WSL2 involves
six non-obvious compatibility issues absent from any published CTranslate2 4.7.1
or PyTorch 2.6/2.7 guide.
Table~\ref{tab:constraints} summarises all six.

\begin{table}[htbp]
\caption{Hardware Constraints and Solutions.}
\label{tab:constraints}
\centering
\resizebox{\columnwidth}{!}{%
\setlength{\tabcolsep}{4pt}%
\begin{tabular}{>{\raggedright\arraybackslash}p{0.22\columnwidth}
               >{\raggedright\arraybackslash}p{0.38\columnwidth}
               >{\raggedright\arraybackslash}p{0.27\columnwidth}}
\toprule
Constraint & Root Cause & Solution \\
\midrule
CT2 INT8 fails & cuBLAS cu124 lacks sm\,120 INT8 tensor core support &
  Use float16; auto-probe \\
\addlinespace
Corrupt pip install & Concurrent pip processes corrupt shared objects &
  Single sequential install \\
\addlinespace
Tokenizer deadlock & Fast tokenizer thread pool forks into DataLoader workers;
  threads do not restart &
  Set TOKENIZERS\_PARALLELISM=false \\
\addlinespace
prefetch error & Trainer sets prefetch=2 with num\_workers=0;
  PyTorch rejects non-None value &
  prefetch=None with 0 workers \\
\addlinespace
fp16 GradScaler & half() weights with fp16=True raises
  ValueError on gradient unscale &
  Keep float32; use bf16=True \\
\addlinespace
CT2 batch OOM & 900+ sequences at once exceeds CUDA memory &
  max\_batch\_size=32 \\
\bottomrule
\end{tabular}%
}
\end{table}

\textbf{Constraint 1 --- CTranslate2 INT8 on sm\_120.}
CTranslate2 INT8 calls cuBLAS INT8 tensor core operations~\cite{dettmers2022llmint8}.
On CUDA~12.4 with Blackwell sm\_120, these raise \texttt{CUBLAS\_STATUS\_NOT\_SUPPORTED}.
Float16 uses standard CUDA matmul paths supported via PTX JIT.
An active probe translates a 20-token Bengali sentence to select compute type at
load time; probes shorter than 10 tokens do not trigger INT8 operations and produce
false positives.

\textbf{Constraint 2 --- Concurrent pip install SIGBUS.}
An initial diagnosis attributed SIGBUS during \texttt{import torch} to AMD Ryzen AMX
incompatibility.
This was incorrect: AMD Ryzen~5 240 has no \texttt{amx\_bf16}, \texttt{amx\_tile},
or \texttt{amx\_int8} flags in \texttt{/proc/cpuinfo}.
Overlapping writes from two concurrent \texttt{pip install} processes produced a
truncated \texttt{.so} shared object, causing a bus error at dynamic link time.
A clean single-process reinstall eliminates the SIGBUS.

\textbf{Constraint 3 --- HuggingFace Tokenizers fork deadlock.}
NLLB-200's fast tokenizer maintains an internal Rust thread pool.
Forking the main process for DataLoader workers inherits pool state but does not
restart threads; subsequent tokenization in a worker blocks indefinitely.
Setting \texttt{TOKENIZERS\_PARALLELISM=false} before Trainer construction disables
the pool in child processes.

\textbf{Constraint 4 --- DataLoader \texttt{prefetch\_factor} with \texttt{num\_workers=0}.}
PyTorch rejects any non-\texttt{None} \texttt{prefetch\_factor} when
\texttt{num\_workers=0}; HuggingFace Trainer defaults to \texttt{prefetch\_factor=2}.
The fix is to explicitly set \texttt{prefetch\_factor=None} for single-process data
loading.

\textbf{Constraint 5 --- fp16 model weights + fp16 AMP GradScaler conflict.}
Loading model parameters in float16 via \texttt{.half()} and then setting
\texttt{fp16=True} in \texttt{Seq2SeqTrainingArguments} raises
\texttt{ValueError: Attempting to unscale FP16 gradients}.
The fix is to keep parameters in float32 and set \texttt{bf16=True};
BrainFloat16~\cite{micikevicius2018mixed} needs no GradScaler and the
RTX~5050~\cite{nvidia_blackwell} natively supports bf16 tensor cores.

\textbf{Constraint 6 --- CTranslate2 \texttt{translate\_batch} OOM with large inputs.}
Passing 900+ tokenized sequences to \texttt{translate\_batch()} allocates GPU memory
proportional to the full batch before decoding, causing CUDA OOM on 8\,GB VRAM.
Passing \texttt{max\_batch\_size=32} enables internal batching in the C++ layer,
capping peak VRAM allocation.

% ============================================================
\section{Limitations}
\label{sec:limits}
% ============================================================

\textbf{Open-domain vs.\ in-domain BLEU gap.}
The 65.2 BLEU on the curated 90-sentence corpus contrasts with an effective
open-domain score of approximately 28--31 on FLORES-200-equivalent data.
The sacreBLEU figure of 0.15 on Samanantar is a measurement artifact from
tokenization mismatch (Section~\ref{sec:finetune}), not a quality regression.
Both figures are internally consistent with published NLLB-200 behavior.
The 36.7-point gap between in-domain and open-domain reflects genuine domain
shift between the curated corpus domains (Health, Science, Literature) and the
web-crawled Samanantar distribution.

\textbf{VRAM ceiling for large models.}
Any model exceeding approximately 3.5\,GB float16 (e.g., NLLB-3.3B at 13\,GB,
MADLAD-400-3B at 6.5\,GB) cannot be loaded concurrently with the Ollama polish
pass (4.7\,GB); sequential VRAM swapping is required.
IndicTrans2-1B (3.2\,GB) fits within the 8\,GB budget alongside OS overhead.

\textbf{No chrF or COMET evaluation.}
Only BLEU is reported; chrF~\cite{popovic2015chrf} and COMET~\cite{rei2020comet}
would provide more signal on morphologically complex outputs.
These are planned as future work.

\textbf{Dialect and code-switching.}
Informal Bengali text mixes Bengali script and Roman-script English within the
same sentence.
Standard NMT trained on formal parallel text handles these mixed-script tokens
poorly; transliterating Roman fragments into Bengali is a common error mode in
NLLB-200 on social media inputs.
No code-switching handling is currently implemented.

% ============================================================
\section{Conclusion}
\label{sec:conclusion}
% ============================================================

Fully offline, privacy-preserving Bengali-to-English NMT at 65.2 BLEU and
97--100\,chars/s is achievable on a \$300 consumer laptop GPU (NVIDIA RTX~5050,
Blackwell sm\_120, 8\,GB GDDR7).
CTranslate2 float16 inference reaches 82--84\% average GPU utilisation with a
1.9\,GB model footprint, leaving 6\,GB free for concurrent system use.
LoRA fine-tuning of the same 600M-parameter model completes in \TRAINDURATION{}
on that hardware with evaluation loss decreasing monotonically across all three
epochs.
Neither result required cloud infrastructure, API keys, or multi-GPU scaling.

The six Blackwell deployment constraints documented here---cuBLAS INT8 failure,
corrected SIGBUS root cause (concurrent pip install, not AMD AMX),
fast-tokenizer fork-deadlock, \texttt{prefetch\_factor} conflict,
fp16/bf16 GradScaler interaction, and CTranslate2 batch OOM---are absent from
any published guide for CTranslate2 4.7.1 or PyTorch 2.6/2.7.
Practitioners targeting RTX~50-series hardware will encounter at least four of these
six issues; the runtime compute-type probe, \texttt{max\_batch\_size=32} guard, and
bf16 training configuration documented here resolve each reproducibly.

Per-domain analysis reveals that Health scores 80.6 BLEU and News scores 4.7 BLEU
within the same model run---a 75.9-point spread that single-number benchmarks
cannot represent.
This variance, combined with the SacreBLEU tokenization sensitivity demonstrated
here, argues for per-domain evaluation as the minimum standard for reporting
Bengali NMT results.

Future directions include upgrading to IndicTrans2-1B (+8--11 BLEU on FLORES-200
within the same VRAM budget), scaling the LoRA corpus to the full 9M Samanantar
pairs, and adding chrF and COMET evaluation.
The 186-test TDD suite, resource monitoring subsystem, and hardware constraint
documentation are published at \url{https://github.com/sbisw/translate} as an
open-source reference for consumer-GPU NMT deployment on Blackwell hardware.

% ============================================================
\section*{Acknowledgments}
% ============================================================

The author thanks the teams at Meta AI (NLLB-200), AI4Bharat (IndicTrans2),
HuggingFace (Transformers, PEFT, Datasets), and the OpenNMT project (CTranslate2)
for releasing their models and tools under open-source licenses.

% ============================================================
\bibliographystyle{IEEEtran}
\begin{thebibliography}{99}

\bibitem{nllb2022}
NLLB Team, ``No Language Left Behind: Scaling Human-Centered Machine Translation,''
\textit{arXiv preprint arXiv:2207.04672}, Jul. 2022.
[Online]. Available: \url{https://arxiv.org/abs/2207.04672}

\bibitem{gala2023indictrans2}
J. Gala et al., ``IndicTrans2: Towards High-Quality and Accessible Machine Translation
Models for all 22 Scheduled Indian Languages,''
\textit{Transactions on Machine Learning Research}, 2023.
[Online]. Available: \url{https://arxiv.org/abs/2305.16307}

\bibitem{hu2021lora}
E. J. Hu, Y. Shen, P. Wallis, Z. Allen-Zhu, Y. Li, S. Wang, and W. Chen,
``LoRA: Low-Rank Adaptation of Large Language Models,''
in \textit{Proc. ICLR 2022}, Apr. 2022.
[Online]. Available: \url{https://arxiv.org/abs/2106.09685}

\bibitem{samanantar2022}
G. Ramesh et al., ``Samanantar: The Largest Publicly Available Parallel Corpora
Collection for 11 Indic Languages,''
\textit{Trans. Assoc. Comput. Linguistics}, vol.~10, pp.~145--162, 2022.
DOI: 10.1162/tacl\_a\_00452

\bibitem{vaswani2017attention}
A. Vaswani, N. Shazeer, N. Parmar, J. Uszkoreit, L. Jones, A. N. Gomez,
\L. Kaiser, and I. Polosukhin, ``Attention Is All You Need,''
in \textit{Proc. NeurIPS 2017}, pp.~5998--6008, 2017.

\bibitem{ctranslate2_klein}
G. Klein, Y. Kim, Y. Deng, J. Senellart, and A. Rush,
``OpenNMT: Open-Source Toolkit for Neural Machine Translation,''
in \textit{Proc. ACL 2017, System Demonstrations}, pp.~67--72, 2017.
(CTranslate2 is the successor inference engine.)
[Online]. Available: \url{https://github.com/OpenNMT/CTranslate2}

\bibitem{peft2022}
S. Mangrulkar, S. Ghosh, L. Wang, J. Guo, and B. M. Saqur,
``PEFT: State-of-the-art Parameter-Efficient Fine-Tuning,''
HuggingFace, 2022.
[Online]. Available: \url{https://github.com/huggingface/peft}

\bibitem{wolf2020transformers}
T. Wolf et al., ``Transformers: State-of-the-Art Natural Language Processing,''
in \textit{Proc. EMNLP 2020: System Demonstrations}, pp.~38--45, 2020.

\bibitem{papineni2002bleu}
K. Papineni, S. Roukos, T. Ward, and W.-J. Zhu,
``BLEU: A Method for Automatic Evaluation of Machine Translation,''
in \textit{Proc. ACL 2002}, pp.~311--318, 2002.

\bibitem{post2018call}
M. Post, ``A Call for Clarity in Reporting BLEU Scores,''
in \textit{Proc. WMT 2018}, pp.~186--191, 2018.

\bibitem{popovic2015chrf}
M. Popovi\'{c}, ``chrF: character n-gram F-score for automatic MT evaluation,''
in \textit{Proc. WMT 2015}, pp.~392--395, 2015.

\bibitem{callison2006re}
C. Callison-Burch, M. Osborne, and P. Koehn,
``Re-evaluating the Role of BLEU in Machine Translation Research,''
in \textit{Proc. EACL 2006}, pp.~249--256, 2006.

\bibitem{liu2020multilingual}
Y. Liu et al., ``Multilingual Denoising Pre-training for Neural Machine Translation,''
\textit{Trans. Assoc. Comput. Linguistics}, vol.~8, pp.~726--742, 2020.

\bibitem{fan2021beyond}
A. Fan et al., ``Beyond English-Centric Multilingual Machine Translation,''
\textit{J. Machine Learning Research}, vol.~22, no.~107, pp.~1--48, 2021.

\bibitem{li2021prefix}
X. L. Li and P. Liang, ``Prefix-Tuning: Optimizing Continuous Prompts for Generation,''
in \textit{Proc. ACL 2021}, pp.~4582--4597, 2021.

\bibitem{lester2021power}
B. Lester, R. Al-Rfou, and N. Constant,
``The Power of Scale for Parameter-Efficient Prompt Tuning,''
in \textit{Proc. EMNLP 2021}, pp.~3045--3059, 2021.

\bibitem{arivazhagan2019massively}
N. Arivazhagan et al., ``Massively Multilingual Neural Machine Translation in the Wild:
Findings and Challenges,''
\textit{arXiv preprint arXiv:1907.05019}, 2019.

\bibitem{llamacpp}
G. Gerganov, ``llama.cpp: Port of Facebook's LLaMA model in C/C++,'' 2023.
[Online]. Available: \url{https://github.com/ggerganov/llama.cpp}

\bibitem{hasan2020xl}
T. Hasan et al., ``XL-Sum: Large-Scale Multilingual Abstractive Summarization
for 44 Languages,''
in \textit{Findings of ACL-IJCNLP 2021}, pp.~4693--4703, 2021.

\bibitem{fadaee2017data}
M. Fadaee, A. Bisazza, and C. Monz,
``Data Augmentation for Low-Resource Neural Machine Translation,''
in \textit{Proc. ACL 2017}, pp.~567--573, 2017.

\bibitem{wu2016googles}
Y. Wu et al., ``Google's Neural Machine Translation System: Bridging the Gap between
Human and Machine Translation,''
\textit{arXiv preprint arXiv:1609.08144}, 2016.

\bibitem{ethnologue2023}
D. M. Eberhard, G. F. Simons, and C. D. Fennig (Eds.),
\textit{Ethnologue: Languages of the World}, 26th ed.,
SIL International, Dallas, TX, 2023.
[Online]. Available: \url{https://www.ethnologue.com}

\bibitem{shore2004fail}
J. Shore, ``Fail Fast,''
\textit{IEEE Software}, vol.~21, no.~5, pp.~21--25, Sep./Oct. 2004.

\bibitem{micikevicius2018mixed}
P. Micikevicius et al., ``Mixed Precision Training,''
in \textit{Proc. ICLR 2018}, 2018.
[Online]. Available: \url{https://arxiv.org/abs/1710.03740}

\bibitem{dettmers2022llmint8}
T. Dettmers, M. Lewis, Y. Belkada, and L. Zettlemoyer,
``LLM.int8(): 8-bit Matrix Multiplication for Transformers at Scale,''
in \textit{Proc. NeurIPS 2022}, 2022.
[Online]. Available: \url{https://arxiv.org/abs/2208.07339}

\bibitem{pytorch2024}
A. Paszke et al., ``PyTorch: An Imperative Style, High-Performance Deep Learning Library,''
in \textit{Proc. NeurIPS 2019}, pp.~8024--8035, 2019.
[Online]. Available: \url{https://arxiv.org/abs/1912.01703}

\bibitem{nvidia_blackwell}
NVIDIA Corporation, ``NVIDIA Blackwell Architecture Technical Brief,''
Technical Report, 2024.
[Online]. Available: \url{https://resources.nvidia.com/en-us-blackwell-architecture}

\bibitem{kudugunta2023madlad}
S.~Kudugunta et~al.,
``MADLAD-400: A Multilingual And Document-Level Large Audited Dataset,''
\textit{arXiv preprint arXiv:2309.04662}, 2023.

\bibitem{seamlessm4t2023}
Meta AI,
``SeamlessM4T---Massively Multilingual \& Multimodal Machine Translation,''
\textit{arXiv preprint arXiv:2308.11596}, 2023.

\bibitem{rei2020comet}
R.~Rei, C.~Stewart, A.~C.~Farinha, and A.~Lavie,
``COMET: A Neural Framework for MT Evaluation,''
in \textit{Proc. EMNLP 2020}, pp.~2685--2702, 2020.

\end{thebibliography}

\end{document}
